\begin{proposition}[exact global-section criterion]
\label{prop:criterion}
Let $A$ be a line bundle on $S$.  If
\[
 H^0(S,\widetilde\Omega^1_S\otimes A)=0,
 \tag{2.6}
\]
then the inclusion in (2.5) induces an isomorphism
\[
 H^0(S,\Omega^1_X|_S\otimes A)
 \simeq H^0(S,\mathscr K_S\otimes A).
 \tag{2.7}
\]
If in addition
\[
 H^0(S,N^*_{S/X}\otimes A)=0,
 \tag{2.8}
\]
then there is a canonical isomorphism
\[
 H^0(S,\Omega^1_X|_S\otimes A)
 \simeq \ker\!\left(
 H^0(S,\mathscr T_S\otimes A)
 \xrightarrow{\;\delta_A\;}
 H^1(S,N^*_{S/X}\otimes A)
 \right),
 \tag{2.9}
\]
where $\delta_A$ is the connecting map of (2.4) tensored by $A$.
Consequently, under (2.6) and (2.8), the ambient group vanishes if and only
if $\delta_A$ is injective.
\end{proposition}

\begin{proof}
Tensor (2.5) by the locally free sheaf $A$ and take global sections.  If
(2.6) holds, the long exact sequence gives (2.7).  Tensor (2.4) by $A$ and
take its first four terms.  Condition (2.8) identifies the middle global
sections with the kernel of the connecting map.  Combining the two
identifications gives (2.9).
\end{proof}

\begin{remark}
The map $\delta_A$ is not an auxiliary hypothesis introduced for
terminological convenience: it is the exact obstruction to lifting a torsion
section of $\Omega^1_S\otimes A$ through the conormal extension.  In
\v{C}ech terms, choose local lifts of a section to
$\mathscr K_S\otimes A$; their pairwise differences are sections of
$N^*_{S/X}\otimes A$ and represent $\delta_A$.  The two endpoint groups in
(2.6) and (2.8) contain no information about this obstruction.
\end{remark}

\section{A conductor-supported section}

The normalization--conductor mechanism in this section is a self-contained
restatement (and slight sheaf-theoretic generalization) of the source
manuscript's Lemma~10.3; it is included so that the independent test family
can be checked without importing an unprinted step, and no priority is
claimed for it.  We use it to produce a nonzero element of the kernel in
(2.9).  Write
\(\theta_g:=i^*(dg)\in H^0(S,\Omega^1_X|_S)\).  Whenever we write
$dg|_S$ below it means this ambient restriction $i^*(dg)$, not the intrinsic
differential $d(g|_S)=0$ of the constant function on the fibre.  Suppose
$g:X\to\Delta$ is a holomorphic map from a smooth complex threefold to a
disc, $S=g^{-1}(0)$ is reduced, and, writing $D=\Sing(S)$, every point of
$D$ has a normal-crossing (self-crossing) map-germ model: in local coordinates
$(u,v,w)$, $g=h\,uv$ with $h$ a holomorphic unit.  After replacing $u$ by
$hu$, we use the equivalent model
\[
 (X,g,S)\simeq\bigl(\C^3,uv,\{uv=0\}\bigr).
 \tag{3.1}
\]
At points of $S\setminus D$ use the smooth model $(\C^3,\{u=0\})$.
Let $\eta:\widetilde S\to S$ be the normalization.
For the application below, $D$ is smooth and $\eta^{-1}(D)$ is the disjoint
union of two copies of $D$.

